\documentclass[scrarticle]{scrartcl}

\usepackage[T1]{fontenc}
\usepackage[utf8]{inputenc}

\usepackage{amsmath}
\usepackage{amssymb}
\usepackage{array}
\usepackage{authblk}
   \renewcommand\Affilfont{\small}
\usepackage[english]{babel}
\usepackage[style=english]{csquotes}
\usepackage{dcolumn}
   \newcolumntype{d}[1]{D{.}{.}{#1}}
\usepackage{enumitem}
\usepackage{float}
\usepackage{graphicx}
   \graphicspath{{./figures/}}
\usepackage[hidelinks]{hyperref}
   \urlstyle{rm}
\usepackage[os=win]{menukeys}
\usepackage[authoryear]{natbib}
\usepackage{pdflscape}
\usepackage{wasysym}

\setcitestyle{aysep={}}

\deffootnote{1.5em}{1em}{\makebox[1.5em][l]{\thefootnotemark}}
   \setlength{\skip\footins}{1.5em}
   \setlength{\footnotesep}{1em}

\addto\extrasenglish{
   \def\sectionautorefname{Section}
   \def\subsectionautorefname{Subsection}
}

\title{Preregistration}
\subtitle{Experiment 3b:\\Does your upbringing affect your moral responsibility as an adult?\\Follow-up study}

\author[1]{Pascale Willemsen}
\author[2]{Alexander Max Bauer}
\author[1]{Pauline Maresi Kinzler }

\affil[1]{ University of Zurich}
\affil[2]{ Carl von Ossietzky University of Oldenburg}

\date{\small preregistered at \url{https://osf.io/h8nmg/} on April 21, 2026}


\begin{document}
\maketitle


%%%%%%%%%%%%%%%
% DESCRIPTION %
%%%%%%%%%%%%%%%
\clearpage
\section{Description}
This is a follow-up based on Experiment 3a.


%%%%%%%%%%%%%%%%%%%
% DATA COLLECTION %
%%%%%%%%%%%%%%%%%%%
\section{Data Collection}
\begin{itemize}[label=\Square,noitemsep]
   \item Yes, we already collected the data.
   \item[\XBox] No, no data have been collected for this study yet.
   \item It's complicated. We have already collected some data but explain in Question $8$ why readers may consider this a valid pre-registration nevertheless.
\end{itemize}


%%%%%%%%%%%%%%
% HYPOTHESIS %
%%%%%%%%%%%%%%
\section{Hypothesis}
Experiment 3b is identical to Experiment 3a in all respects, except for one thing:
Instead of presenting the two DVs Blame/Praise and True Self on the same page, we now present them on different pages.
Participants see the Blame or Praise question first (depending on the conditions they were assigned to) below the vignette, they answer the question, and proceed to the next screen.
On the next screen, they see the vignette again, followed by the True Self DV.
Participants cannot go back to the previous page to change their answer.

For both pairs of cases (two blame and two praise), we predict a significant difference.

We predict that the ignorant agent is considered less blameworthy than their non-ignorant counterpart.

We predict that the ignorant agent is considered more praiseworthy than their non-ignorant counterpart.


%%%%%%%%%%%%%%%%%%%%%%%
% DEPENDENT VARIABLES %
%%%%%%%%%%%%%%%%%%%%%%%
\section{Dependent Variables}


%%%%%%%%%%%%%%%%%%%%%%%%%%%%%
% DEPENDENT VARIABLES – DV1 %
%%%%%%%%%%%%%%%%%%%%%%%%%%%%%
\subsection*{DV1 (Blame)}
On the scale below, please choose the number that best represents the degree to which you think Tom is morally blameworthy (if at all) for spitting on the black man.

\vspace{0.75em}
\noindent Scale and labels:

\vspace{0.75em}
\begin{tabular}{r @{ = } l}
    $1$   & not at all blameworthy   \\
    $4$   & somewhat blameworthy     \\
    $7$   & completely blameworthy   \\
\end{tabular}


%%%%%%%%%%%%%%%%%%%%%%%%%%%%%
% DEPENDENT VARIABLES – DV2 %
%%%%%%%%%%%%%%%%%%%%%%%%%%%%%
\subsection*{DV2 (Praise)}
On the scale below, please choose the number that best represents the degree to which you think Tom is morally praiseworthy (if at all) for helping the man up.

\vspace{0.75em}
\noindent Scale and labels:

\vspace{0.75em}
\begin{tabular}{r @{ = } l}
    $1$   & not at all praiseworthy   \\
    $4$   & somewhat praiseworthy     \\
    $7$   & completely praiseworthy   \\
\end{tabular}


%%%%%%%%%%%%%%%%%%%%%%%%%%%%%
% DEPENDENT VARIABLES – DV3 %
%%%%%%%%%%%%%%%%%%%%%%%%%%%%%
\subsection*{DV3 (True Self)}
One the scale below, please circle the number that best represents the extent to which you agree that what Tom did expressed his true self---the person he really is deep down inside.


\vspace{0.75em}
\noindent Scale and labels:

\vspace{0.75em}
\begin{tabular}{r @{ = } l}
    $1$   & disagree completely          \\
    $4$   & neither agree nor disagree   \\
    $7$   & agree completely             \\
\end{tabular}


%%%%%%%%%%%%%%
% CONDITIONS %
%%%%%%%%%%%%%%
\section{Conditions}
This is a $2 \times 2$ between-subjects design with the independent variables Ignorance (Ignorant, Non-Ignorant) and Moral Valence of the Action (Blameworthy, Praiseworthy).
As stimuli, we used the materials shared by Faraci and Shoemaker.


%%%%%%%%%%%%%%%%%%%
% CONDITIONS – NB %
%%%%%%%%%%%%%%%%%%%
\subsection*{Non-Ignorant, Blameworthy}
Tom is a white male who was raised in New Orleans.
Growing up, he was taught to respect all people equally.
Nevertheless, as an adult, he decided to become a proud racist, someone who believes that all non-white people are inferior and that he has a moral obligation to humiliate them when he gets a chance.
At the age of $25$, Tom moves to another town.
Walking outside his home, he sees a black man who has tripped and fallen.
In keeping with his moral beliefs, Tom spits on the man as he passes by.


%%%%%%%%%%%%%%%%%%%
% CONDITIONS – IB %
%%%%%%%%%%%%%%%%%%%
\subsection*{Ignorant, Blameworthy}
Tom is a white male who was raised in New Orleans.
Growing up, he was taught to respect all people equally.
Nevertheless, as an adult, he decided to become a proud racist, someone who believes that all non-white people are inferior and that he has a moral obligation to humiliate them when he gets a chance.
At the age of $25$, Tom moves to another town.
Walking outside his home, he sees a black man who has tripped and fallen.
In keeping with his moral beliefs, Tom spits on the man as he passes by.


%%%%%%%%%%%%%%%%%%%
% CONDITIONS – NP %
%%%%%%%%%%%%%%%%%%%
\subsection*{Non-Ignorant, Praiseworthy}
Tom is a white male who was raised in New Orleans.
Growing up, he was taught to respect all people equally.
Nevertheless, as an adult, he decided to become a proud racist, someone who believes that all non-white people are inferior and that he has a moral obligation to humiliate them when he gets a chance.
At the age of $25$, Tom moves to another town.
Walking outside his home, he sees a black man trip and fall.
Usually, Tom would spit on the man.
But this time, Tom goes against his current moral beliefs, and helps the man up instead.


%%%%%%%%%%%%%%%%%%%
% CONDITIONS – IP %
%%%%%%%%%%%%%%%%%%%
\subsection*{Ignorant, Praiseworthy}
Tom is a white male who was raised on an isolated island in the bayous of Louisiana.
Growing up, he was taught to believe that all non-white people are inferior and that he has a moral obligation to humiliate them when he gets a chance.
As an adult, he decided to become a proud racist, embracing what he was taught.
At the age of $25$, Tom moves to another town.
Walking outside his home, he sees a black man trip and fall.
Usually, Tom would spit on the man.
But this time, Tom goes against his current moral beliefs, and helps the man up instead.


%%%%%%%%%%%%
% ANALYSES %
%%%%%%%%%%%%
\section{Analyses}
Just as Faraci and Shoemaker ($2014$), we conduct two t-tests, comparing the two blame and the two praise conditions with one another.


%%%%%%%%%%%%
% OUTLIERS %
%%%%%%%%%%%%
\section{Outliers and Exclusions}
Just as in Experiments $1$, $2$, and $3$a, we use Prolific's screeners and only allow participation if the following conditions are met:

\begin{itemize}
   \item[] First Language: English
   \item[] Country of Residence: US
\end{itemize}

\noindent We use a screener to exclude participants from other studies related to this project.


%%%%%%%%%%%%%%%
% SAMPLE SIZE %
%%%%%%%%%%%%%%%
\section{Sample Size}
An \textit{a priori} power analysis was conducted to determine the required sample size for each of the two two-tailed independent samples t-test.
Based on Faraci and Shoemaker's earlier work, a large effect size ($d \geq 0.80$) was assumed.
With an alpha level of $\alpha = 0.01$ and a planned sample size of $N = 280$ ($n = 70$ per group), the achieved statistical power ($1 - \beta$) is approximately $0.935$ ($93.5\,\%$).
This sample size ensures that the study is sufficiently powered to detect the hypothesized effect while maintaining a strict control for Type I errors.


%%%%%%%%%
% OTHER %
%%%%%%%%%
\section{Other}


%%%%%%%%%%%%%%%%%%
% OTHER – PRAISE %
%%%%%%%%%%%%%%%%%%
\subsection*{Praise Explanation}
In this study, we are particularly interested in why people like you consider others blame- or praiseworthy.
Please explain in as many words as you like why you considered Tom praiseworthy to the degree you chose.


%%%%%%%%%%%%%%%%%
% OTHER – BLAME %
%%%%%%%%%%%%%%%%%
\subsection*{Blame Explanation}
In this study, we are particularly interested in why people like you consider others blame- or praiseworthy.
Please explain in as many words as you like why you considered Tom blameworthy to the degree you chose.


%%%%%%%%%%%%%%%%%%%
% OTHER – CONTROL %
%%%%%%%%%%%%%%%%%%%
\subsection*{Control Variables}
To ensure that our results are robust to individual differences in philosophical intuitions, we include control variables in our ANOVA (ANCOVA), including Gender, Age, Education, Personal Experience, Free Will, Social Determinism, Political Orientation, and Metaethical Stance.


%%%%%%%%
% NAME %
%%%%%%%%
\section{Name}
Replication of Faraci and Shoemaker $2014$ with page breaks between DVs


%%%%%%%%
% TYPE %
%%%%%%%%
\section{Type of Project}
For record keeping purposes, please tell us the type of study you are pre-registering.

\begin{itemize}[label=\Square,noitemsep]
   \item Class project or assignment
   \item[\XBox] Experiment
   \item Survey
   \item Observational/archival study
   \item Other (describe below)
\end{itemize}


\end{document}
